\documentclass[11pt]{article}
\usepackage{amsmath,amssymb,amsthm,mathtools}
\usepackage[a4paper,margin=1in]{geometry}
\usepackage{enumitem}
\usepackage{hyperref}
\usepackage{tikz-cd}

\title{Concrete de Rham Cohomology on \(S^2\) and \(T^2\),\\
	Winding and Area Forms, Euler Class, and Gauss--Bonnet}
\author{}
\date{}

\newtheorem{theorem}{Theorem}
\newtheorem{proposition}{Proposition}
\newtheorem{remark}{Remark}
\newtheorem{exercise}{Exercise}

\begin{document}
	\maketitle
	
	\tableofcontents
	
	\section{Goal}
	
	We want a compute-first understanding of:
	\begin{itemize}[leftmargin=2em]
		\item de Rham cohomology through explicit differential forms,
		\item the winding form on \(\mathbb{R}^2\setminus\{0\}\),
		\item the area form on \(S^2\),
		\item the basic forms \(dx,dy,dx\wedge dy\) on \(T^2\),
		\item the overlap/gluing viewpoint via \v{C}ech cocycles and transition functions,
		\item the Euler class of \(TS^2\),
		\item and Gauss--Bonnet.
	\end{itemize}
	
	The central theme is:
	\[
	\boxed{\text{closed} \neq \text{exact}}
	\]
	and this failure is measured by topology.
	
	\section{de Rham cohomology}
	
	For a smooth manifold \(M\),
	\[
	H^k_{\mathrm{dR}}(M)
	=
	\frac{\{\text{closed }k\text{-forms}\}}{\{\text{exact }k\text{-forms}\}}.
	\]
	
	So:
	\begin{itemize}[leftmargin=2em]
		\item \(\omega\) is \emph{closed} if \(d\omega=0\),
		\item \(\omega\) is \emph{exact} if \(\omega=d\eta\),
		\item two closed forms define the same cohomology class if they differ by an exact form.
	\end{itemize}
	
	Thus a nonzero cohomology class is represented by a closed form which is not exact.
	
	\section{Warm-up: the winding form on \(\mathbb{R}^2\setminus\{0\}\)}
	
	Consider
	\[
	\eta=\frac{-y\,dx+x\,dy}{x^2+y^2}.
	\]
	
	\subsection{Polar-coordinate computation}
	
	Let
	\[
	x=r\cos\theta,\qquad y=r\sin\theta.
	\]
	Then
	\[
	dx=\cos\theta\,dr-r\sin\theta\,d\theta,
	\qquad
	dy=\sin\theta\,dr+r\cos\theta\,d\theta.
	\]
	Hence
	\begin{align*}
		-y\,dx+x\,dy
		&=
		-r\sin\theta(\cos\theta\,dr-r\sin\theta\,d\theta)
		+r\cos\theta(\sin\theta\,dr+r\cos\theta\,d\theta)\\
		&=
		-r\sin\theta\cos\theta\,dr+r^2\sin^2\theta\,d\theta
		+r\sin\theta\cos\theta\,dr+r^2\cos^2\theta\,d\theta\\
		&=
		r^2\,d\theta.
	\end{align*}
	Since \(x^2+y^2=r^2\), we get
	\[
	\eta=d\theta.
	\]
	
	So \(\eta\) is locally exact.
	
	\subsection{Closed but not exact}
	
	Since \(\eta=d\theta\) locally,
	\[
	d\eta=0.
	\]
	Thus \(\eta\) is closed on \(\mathbb{R}^2\setminus\{0\}\).
	
	Now integrate around the unit circle \(\gamma(t)=(\cos t,\sin t)\), \(0\le t\le 2\pi\):
	\[
	\int_\gamma \eta = \int_0^{2\pi} d\theta = 2\pi.
	\]
	If \(\eta=df\) globally, then \(\int_\gamma \eta=0\), contradiction. So \(\eta\) is not exact.
	
	Therefore
	\[
	[\eta]\neq 0 \in H^1_{\mathrm{dR}}(\mathbb{R}^2\setminus\{0\}).
	\]
	
	\begin{remark}
		This is the prototype: there are local potentials, but no global single-valued potential.
	\end{remark}
	
	\section{The sphere \(S^2\): area form and cohomology}
	
	Let \(S^2\subset \mathbb{R}^3\) be the unit sphere. Use spherical coordinates
	\[
	X(\theta,\phi)=
	(\sin\phi\cos\theta,\sin\phi\sin\theta,\cos\phi),
	\qquad
	0<\theta<2\pi,\ 0<\phi<\pi.
	\]
	
	\subsection{Metric and area form}
	
	A standard computation gives
	\[
	ds^2=d\phi^2+\sin^2\phi\,d\theta^2.
	\]
	Hence the area form is
	\[
	dA=\sin\phi\,d\theta\wedge d\phi.
	\]
	
	\subsection{Closed but not exact}
	
	Since \(S^2\) is \(2\)-dimensional, every \(3\)-form is zero, so
	\[
	d(dA)=0.
	\]
	Thus \(dA\) is closed.
	
	Now compute its integral:
	\begin{align*}
		\int_{S^2} dA
		&=
		\int_0^{2\pi}\int_0^\pi \sin\phi\,d\phi\,d\theta\\
		&=
		\int_0^{2\pi}\left[-\cos\phi\right]_{\phi=0}^{\phi=\pi}\,d\theta\\
		&=
		\int_0^{2\pi}2\,d\theta\\
		&=4\pi.
	\end{align*}
	If \(dA=d\alpha\) globally, then by Stokes,
	\[
	\int_{S^2}dA=\int_{S^2}d\alpha=\int_{\partial S^2}\alpha=0,
	\]
	contradiction. So \(dA\) is not exact.
	
	Therefore
	\[
	[dA]\neq 0\in H^2_{\mathrm{dR}}(S^2).
	\]
	
	\subsection{Why \(H^1_{\mathrm{dR}}(S^2)=0\)}
	
	Concretely, every loop on \(S^2\) bounds a region. So if \(\omega\) is a closed \(1\)-form and \(\gamma=\partial D\), then
	\[
	\int_\gamma \omega=\int_D d\omega=0.
	\]
	Thus all periods vanish, so there is no winding-type obstruction. Hence every closed \(1\)-form is exact:
	\[
	H^1_{\mathrm{dR}}(S^2)=0.
	\]
	
	So:
	\[
	H^0_{\mathrm{dR}}(S^2)\cong \mathbb{R},
	\qquad
	H^1_{\mathrm{dR}}(S^2)=0,
	\qquad
	H^2_{\mathrm{dR}}(S^2)\cong \mathbb{R},
	\]
	with \([dA]\) generating \(H^2\).
	
	\section{Local primitives on \(S^2\) and overlap winding}
	
	Cover \(S^2\) by
	\[
	U_N=S^2\setminus\{N\},\qquad U_S=S^2\setminus\{S\}.
	\]
	
	On \(U_N\), define
	\[
	\alpha_N=(1-\cos\phi)\,d\theta.
	\]
	Then
	\[
	d\alpha_N=d(1-\cos\phi)\wedge d\theta
	=\sin\phi\,d\phi\wedge d\theta
	=-\sin\phi\,d\theta\wedge d\phi.
	\]
	So up to orientation convention,
	\[
	d\alpha_N=dA.
	\]
	
	On \(U_S\), define
	\[
	\alpha_S=-(1+\cos\phi)\,d\theta.
	\]
	Again,
	\[
	d\alpha_S=dA
	\]
	up to the same convention.
	
	On the overlap \(U_N\cap U_S\),
	\[
	\alpha_N-\alpha_S=2\,d\theta.
	\]
	
	Thus the area form is locally exact on each chart, but the local primitives differ by the winding form \(2\,d\theta\) on the overlap.
	
	This is the concrete obstruction to global exactness.
	
	\section{The torus \(T^2=\mathbb{R}^2/\mathbb{Z}^2\)}
	
	Let \(T^2=\mathbb{R}^2/\mathbb{Z}^2\), with coordinates \((x,y)\) modulo integers.
	
	\subsection{Closed \(1\)-forms \(dx,dy\)}
	
	The forms \(dx\) and \(dy\) descend to globally defined \(1\)-forms on \(T^2\), and
	\[
	d(dx)=0,\qquad d(dy)=0.
	\]
	So both are closed.
	
	Now consider the loops
	\[
	\gamma_x(t)=(t,0)\!\!\mod \mathbb{Z}^2,\qquad
	\gamma_y(t)=(0,t)\!\!\mod \mathbb{Z}^2,
	\qquad 0\le t\le 1.
	\]
	Then
	\[
	\int_{\gamma_x}dx=1,\qquad \int_{\gamma_x}dy=0,
	\]
	\[
	\int_{\gamma_y}dx=0,\qquad \int_{\gamma_y}dy=1.
	\]
	
	So neither \(dx\) nor \(dy\) is exact. Hence
	\[
	[dx],[dy]\neq 0 \in H^1_{\mathrm{dR}}(T^2),
	\]
	and they are linearly independent.
	
	Therefore
	\[
	H^1_{\mathrm{dR}}(T^2)\cong \mathbb{R}^2
	\]
	with basis \([dx],[dy]\).
	
	\subsection{The top-degree class \(dx\wedge dy\)}
	
	Consider
	\[
	dx\wedge dy.
	\]
	It is closed because
	\[
	d(dx\wedge dy)=0.
	\]
	Its integral is
	\[
	\int_{T^2}dx\wedge dy=\int_0^1\int_0^1 dx\,dy=1.
	\]
	So it cannot be exact.
	
	Hence
	\[
	[dx\wedge dy]\neq 0 \in H^2_{\mathrm{dR}}(T^2),
	\]
	and it generates \(H^2_{\mathrm{dR}}(T^2)\cong \mathbb{R}\).
	
	Thus
	\[
	H^0_{\mathrm{dR}}(T^2)\cong \mathbb{R},
	\qquad
	H^1_{\mathrm{dR}}(T^2)\cong \mathbb{R}^2,
	\qquad
	H^2_{\mathrm{dR}}(T^2)\cong \mathbb{R}.
	\]
	
	\subsection{Ring structure}
	
	The wedge product gives
	\[
	[dx]\wedge[dx]=0,\qquad [dy]\wedge[dy]=0,
	\]
	\[
	[dx]\wedge[dy]=[dx\wedge dy],
	\qquad
	[dy]\wedge[dx]=-[dx\wedge dy].
	\]
	
	So the top class on the torus is the wedge of the two basic winding classes.
	
	\section{Local-to-global viewpoint: \v{C}ech cocycles and line bundles}
	
	Suppose \(\omega\) is a closed \(2\)-form on a manifold \(M\). By the Poincar\'e lemma, on each sufficiently small open set \(U_i\) we can write
	\[
	\omega|_{U_i}=d\alpha_i.
	\]
	On an overlap \(U_i\cap U_j\),
	\[
	d(\alpha_i-\alpha_j)=0.
	\]
	If the overlap is simply connected, then
	\[
	\alpha_i-\alpha_j=df_{ij}
	\]
	for some function \(f_{ij}\).
	
	So local primitives differ by exact terms on overlaps:
	\[
	\alpha_i-\alpha_j=df_{ij}.
	\]
	
	These \(f_{ij}\) measure the obstruction to gluing the local primitives \(\alpha_i\) into one global \(1\)-form.
	
	Exponentiating gives transition functions
	\[
	g_{ij}=e^{if_{ij}},
	\]
	which satisfy cocycle conditions on triple overlaps and define a line bundle.
	
	\begin{remark}
		This is the bridge:
		\[
		\text{closed form}
		\longleftrightarrow
		\text{local primitives}
		\longleftrightarrow
		\text{overlap functions}
		\longleftrightarrow
		\text{transition functions of a bundle.}
		\]
	\end{remark}
	
	\section{The Euler class and the curvature form}
	
	Let \(M\) be an oriented Riemannian surface. On each sufficiently small open set \(U_i\), choose a positively oriented orthonormal frame
	\[
	(e_1^{(i)},e_2^{(i)}).
	\]
	The Levi-Civita connection is described by a local \(1\)-form \(\omega_i\) such that
	\[
	\nabla e_1^{(i)}=\omega_i\,e_2^{(i)},
	\qquad
	\nabla e_2^{(i)}=-\omega_i\,e_1^{(i)}.
	\]
	
	The structure equation says
	\[
	d\omega_i=K\,dA
	\]
	on each chart. So the curvature form \(K\,dA\) is locally exact.
	
	On an overlap \(U_i\cap U_j\), the two oriented orthonormal frames differ by a rotation:
	\[
	\begin{pmatrix}
		e_1^{(j)}\\
		e_2^{(j)}
	\end{pmatrix}
	=
	\begin{pmatrix}
		\cos\theta_{ij} & -\sin\theta_{ij}\\
		\sin\theta_{ij} & \cos\theta_{ij}
	\end{pmatrix}
	\begin{pmatrix}
		e_1^{(i)}\\
		e_2^{(i)}
	\end{pmatrix}.
	\]
	Then
	\[
	\omega_j=\omega_i+d\theta_{ij}.
	\]
	
	So the connection \(1\)-forms do not globally glue, but their differentials do:
	\[
	d\omega_j=d\omega_i=K\,dA.
	\]
	
	On triple overlaps,
	\[
	\theta_{ij}+\theta_{jk}+\theta_{ki}\in 2\pi\mathbb{Z}.
	\]
	Thus
	\[
	n_{ijk}=\frac{1}{2\pi}\bigl(\theta_{ij}+\theta_{jk}+\theta_{ki}\bigr)\in \mathbb{Z}
	\]
	defines an integer \v{C}ech \(2\)-cocycle. This is the Euler class
	\[
	e(TM)\in H^2(M;\mathbb{Z}).
	\]
	
	Its de Rham image is
	\[
	\boxed{
		\frac{1}{2\pi}[K\,dA]=e(TM)\in H^2_{\mathrm{dR}}(M).
	}
	\]
	
	Integrating over \(M\) gives
	\[
	\frac{1}{2\pi}\int_M K\,dA=\int_M e(TM)=\chi(M),
	\]
	so
	\[
	\boxed{
		\int_M K\,dA=2\pi\chi(M).
	}
	\]
	This is Gauss--Bonnet.
	
	\section{Gauss--Bonnet on \(S^2\) and \(T^2\)}
	
	\subsection{Sphere}
	
	For the unit sphere,
	\[
	K\equiv 1.
	\]
	So
	\[
	K\,dA=dA,
	\]
	and
	\[
	\int_{S^2}K\,dA=\int_{S^2}dA=4\pi.
	\]
	Since
	\[
	\chi(S^2)=2,
	\]
	we get
	\[
	\int_{S^2}K\,dA=2\pi\chi(S^2)=4\pi.
	\]
	
	Cohomologically,
	\[
	[K\,dA]=[dA]\neq 0\in H^2_{\mathrm{dR}}(S^2).
	\]
	
	\subsection{Flat torus}
	
	On the flat torus with metric
	\[
	ds^2=dx^2+dy^2,
	\]
	we have
	\[
	K\equiv 0.
	\]
	Hence
	\[
	K\,dA=0,
	\qquad
	[K\,dA]=0\in H^2_{\mathrm{dR}}(T^2),
	\]
	and
	\[
	\int_{T^2}K\,dA=0.
	\]
	Since
	\[
	\chi(T^2)=0,
	\]
	Gauss--Bonnet gives
	\[
	\int_{T^2}K\,dA=2\pi\chi(T^2)=0.
	\]
	
	\subsection{Important distinction}
	
	The area form \(dA\) on \(T^2\) is not exact, because
	\[
	\int_{T^2} dA = \operatorname{Area}(T^2)>0.
	\]
	So
	\[
	[dA]\neq 0\in H^2_{\mathrm{dR}}(T^2).
	\]
	
	But for the flat torus,
	\[
	K\,dA=0.
	\]
	So Gauss--Bonnet is about \([K\,dA]\), not about \([dA]\) by itself.
	
	\section{Explicit stereographic-coordinate proof that the transition function has degree \(2\)}
	
	We now compute the tangent-bundle transition function of \(S^2\) in stereographic coordinates.
	
	\subsection{Stereographic charts}
	
	Let
	\[
	S^2=\{(X,Y,Z)\in\mathbb{R}^3:X^2+Y^2+Z^2=1\},
	\]
	and let
	\[
	N=(0,0,1),\qquad S=(0,0,-1).
	\]
	
	Define stereographic projection from the north pole:
	\[
	\pi_N(X,Y,Z)=\frac{X+iY}{1-Z}.
	\]
	Define stereographic projection from the south pole:
	\[
	\pi_S(X,Y,Z)=\frac{X-iY}{1+Z}.
	\]
	Write
	\[
	z=\pi_N(p),\qquad w=\pi_S(p).
	\]
	
	\begin{remark}
		The choice of \(X-iY\) in the south chart makes the overlap map holomorphic in the oriented convention.
	\end{remark}
	
	\subsection{Overlap map}
	
	For \(p=(X,Y,Z)\in S^2\setminus\{N,S\}\),
	\[
	z=\frac{X+iY}{1-Z},
	\qquad
	w=\frac{X-iY}{1+Z}.
	\]
	Then
	\[
	zw
	=
	\frac{(X+iY)(X-iY)}{(1-Z)(1+Z)}
	=
	\frac{X^2+Y^2}{1-Z^2}.
	\]
	Since \(X^2+Y^2+Z^2=1\),
	\[
	X^2+Y^2=1-Z^2,
	\]
	so
	\[
	zw=1.
	\]
	Hence
	\[
	w=\frac{1}{z}.
	\]
	
	\subsection{Tangent transition function}
	
	The tangent-bundle transition function is the derivative of the coordinate change:
	\[
	g_{SN}(z)=\frac{dw}{dz}.
	\]
	Since
	\[
	w=\frac{1}{z},
	\]
	we get
	\[
	\frac{dw}{dz}=-\frac{1}{z^2}.
	\]
	
	Now restrict to the equator \(Z=0\). There,
	\[
	z=X+iY=e^{i\theta}.
	\]
	So
	\[
	g_{SN}(e^{i\theta})=-e^{-2i\theta}.
	\]
	
	This already shows the transition winds twice around \(S^1\). Depending on whether one uses north-to-south or south-to-north clutching, and depending on the chosen orientation convention, one gets a degree \(+2\) representative, usually recorded as
	\[
	g(e^{i\theta})=e^{2i\theta}.
	\]
	
	The invariant geometric content is:
	\[
	\boxed{\text{the tangent-bundle clutching map winds twice.}}
	\]
	
	Therefore
	\[
	e(TS^2)[S^2]=2.
	\]
	
	\subsection{Connection with \(\mathcal O(2)\)}
	
	Identifying \(S^2\cong\mathbb{CP}^1\), the transition has degree \(2\), so as a complex line bundle,
	\[
	TS^2\cong \mathcal O(2).
	\]
	Hence
	\[
	c_1(TS^2)=2.
	\]
	For an oriented real rank-\(2\) bundle underlying a complex line bundle,
	\[
	e(TS^2)=c_1(TS^2)=2.
	\]
	
	Thus
	\[
	\frac{1}{2\pi}\int_{S^2}K\,dA=2,
	\qquad
	\int_{S^2}K\,dA=4\pi.
	\]
	
	\section{Big comparison}
	
	\subsection{Sphere \(S^2\)}
	
	\[
	H^0_{\mathrm{dR}}(S^2)\cong\mathbb{R},\qquad
	H^1_{\mathrm{dR}}(S^2)=0,\qquad
	H^2_{\mathrm{dR}}(S^2)\cong\mathbb{R}.
	\]
	A generator of \(H^2\) is \([dA]\). Since \(K=1\),
	\[
	[K\,dA]=[dA]\neq 0.
	\]
	Also,
	\[
	\int_{S^2}K\,dA=4\pi=2\pi\chi(S^2).
	\]
	
	\subsection{Torus \(T^2\)}
	
	\[
	H^0_{\mathrm{dR}}(T^2)\cong\mathbb{R},\qquad
	H^1_{\mathrm{dR}}(T^2)\cong\mathbb{R}^2,\qquad
	H^2_{\mathrm{dR}}(T^2)\cong\mathbb{R}.
	\]
	Generators:
	\[
	[dx],[dy]\in H^1,\qquad [dx\wedge dy]\in H^2.
	\]
	
	For the flat torus,
	\[
	[K\,dA]=0,
	\qquad
	\int_{T^2}K\,dA=0=2\pi\chi(T^2).
	\]
	
	\section{Deep geometric meaning}
	
	There are several equivalent ways to describe the same phenomenon:
	\begin{itemize}[leftmargin=2em]
		\item local potentials exist but global potentials may fail to exist;
		\item local frames exist but global frames may fail to exist;
		\item local connection forms exist but may fail to glue globally;
		\item the obstruction is recorded by cohomology, transition functions, and characteristic classes;
		\item curvature is the differential-form shadow of that obstruction.
	\end{itemize}
	
	For \(S^2\), the obstruction is nontrivial:
	\[
	e(TS^2)=2,\qquad \chi(S^2)=2,\qquad \int_{S^2}K\,dA=4\pi.
	\]
	
	For \(T^2\), the curvature obstruction vanishes:
	\[
	e(TT^2)=0,\qquad \chi(T^2)=0,\qquad \int_{T^2}K\,dA=0.
	\]
	
	\section{Exercises}
	
	\begin{exercise}
		Show directly that
		\[
		d\!\left(\frac{-y\,dx+x\,dy}{x^2+y^2}\right)=0
		\quad\text{on }\mathbb{R}^2\setminus\{0\}.
		\]
	\end{exercise}
	
	\begin{exercise}
		Compute
		\[
		\int_{|z|=1}\frac{-y\,dx+x\,dy}{x^2+y^2}.
		\]
		Conclude that the winding form is not exact.
	\end{exercise}
	
	\begin{exercise}
		Verify that
		\[
		dA=\sin\phi\,d\theta\wedge d\phi
		\]
		on the unit sphere.
	\end{exercise}
	
	\begin{exercise}
		Compute
		\[
		\int_{S^2}dA=4\pi
		\]
		and conclude that \(dA\) is not exact.
	\end{exercise}
	
	\begin{exercise}
		Compute
		\[
		d\bigl((1-\cos\phi)d\theta\bigr),\qquad
		d\bigl(-(1+\cos\phi)d\theta\bigr),
		\]
		and compare the two local primitives on the overlap.
	\end{exercise}
	
	\begin{exercise}
		On \(T^2=\mathbb{R}^2/\mathbb{Z}^2\), compute the periods of \(dx\) and \(dy\) on the two basic cycles and prove that \([dx],[dy]\) are linearly independent.
	\end{exercise}
	
	\begin{exercise}
		Compute
		\[
		\int_{T^2}dx\wedge dy
		\]
		for the square torus and conclude that \(dx\wedge dy\) is not exact.
	\end{exercise}
	
	\begin{exercise}
		Check Gauss--Bonnet directly for the unit sphere and for the flat torus.
	\end{exercise}
	
	\begin{exercise}
		In the stereographic computation, explain carefully why the clutching map may appear as \(-e^{-2i\theta}\) or \(e^{2i\theta}\) depending on convention, while the geometric content remains ``winding number \(2\)''.
	\end{exercise}
	
\end{document}